%% kapa: exam=1 round=22 year=2011 period=1 form=A subject=economics subject_ko=경제원론 seq=27 number=67
%% source: Q-Net 감정평가사 기출문제 (cst00309) artlSeq=5036790 file=extracted/1008/5036790/2011년 제22회 감정평가사 1차시험 기출문제/1교시.hwp
%% conversion: kordoc 4.24.51 -> tools/gen1.py (auto); review=proofread-sonnet; flags=none
\begin{question}{67}
가변생산요소가 하나인 기업의 단기비용곡선에 관한 설명으로 옳지 않은 것은? (단, 평균총비용곡선은 U자 모양, 고정비용 존재, 생산요소가격은 불변)
\begin{choices}
  \item 생산량이 증가함에 따라 한계비용이 증가할 때 한계생산물이 체감한다.
  \item 평균가변비용곡선의 최저점은 평균총비용곡선의 최저점보다 좌측에 위치한다.
  \item 한계비용이 평균총비용보다 작을 때 평균총비용이 상승한다.
  \item 한계비용곡선은 평균총비용곡선의 최저점을 통과한다.
  \item 생산량이 증가함에 따라 평균고정비용이 감소한다.
\end{choices}
\end{question}
\answer{3}
